\documentclass[a4paper, 12pt]{article}

\usepackage{utils}

\renewcommand*{\today}{28 January 2026}

\begin{document}

\hotbox{Géometrie 2}{CM 3}{\today}

\begin{proposition}{}{}
    Soient $\E$ et $\mathcal{F}$ deux espaces affines. Une application $f: \E \to \mathcal{F}$ est affine si et seulement si elle préserve les barycentres.
    
    \vspace{0.5cm}
    Autrement dit, pour tout système de points pondérés \text{$\mathcal{A} = ((A_1, \lambda_1), \dots, (A_n, \lambda_n))$} avec $\sum_{i=1}^n \lambda_i \neq 0$, on a
    $$
    f(bar(\mathcal{A})) = bar(f(\mathcal{A}))
    $$
    où $f(\mathcal{A}) := ((f(A_1), \lambda_1), \dots, (f(A_n), \lambda_n))$.
\end{proposition}

\begin{demonstration}
    \begin{description}
        \item[Sens direct.] %%%%%%%%%%%%%%%%%%%%%%%%
        Supposons $f$ affine, en notant $G = bar(\mathcal{A})$, pour tout $A \in \E$, on a
        \begin{align*}
            f(G) &= f(A + \overrightarrow{AG}) \\
        &= f(A + \sum_{i=1}^n \lambda_i \overrightarrow{A A_i}) \\
        &= f(A) + \vec{f}(\sum_{i=1}^n \lambda_i \overrightarrow{A A_i}) \\
        &= f(A) + \sum_{i=1}^n \lambda_i \vec{f}(\overrightarrow{A A_i}) \\
        &= f(A) + \sum_{i=1}^n \lambda_i \overrightarrow{f(A) f(A_i)} \\
        \end{align*}
        D'où
        $$
        \overrightarrow{f(A) f(G)} = \sum_{i=1}^n \lambda_i \overrightarrow{f(A) f(A_i)}
        $$
        Ainsi $f(G)$ est le barycentre de $((f(A_1), \lambda_1), \dots, (f(A_n), \lambda_n))$.
        
        \item[Sens réciproque.] %%%%%%%%%%%%%%%%%%%%%%
        Supposons que $f$ préserve les barycentres. Montrons que $f$ est affine.

        Soit $A \in \E$ et $A' \in \E$ tels que $A' = f(A)$.

        On fixe un point $O \in \E$ et on définit une application $\funcdef{\vec{f}}{\vec{\E}}{\vec{\mathcal{F}}}{\vec{u} = \overrightarrow{OM}}{\overrightarrow{f(O)f(M)}}$

        Soit $\vec{u} \in \vec{\E}$, il existe un unique $M \in \E$ tel que $\vec{u} = \overrightarrow{OM}$.

        \begin{description}
            \item[Homogénéité.] %%%%%%%%%%%%%%%%%%%%%%%
            Soit $P \in \E$ et $\lambda \in \R$ tels que $\overrightarrow{OP} = \lambda \vec{u} = \lambda \overrightarrow{OM}$.
            
            On peut voir $P$ comme un barycentre de $(O, \alpha), (M, \beta)$, on peut trouver $\alpha$ et $\beta$ avec
            \begin{align*}
                \overrightarrow{OP} &= \dfrac{1}{\alpha + \beta} (\alpha \overrightarrow{OO} + \beta \overrightarrow{OM}) \\
                &= \dfrac{\beta}{\alpha + \beta} \overrightarrow{OM} \implies \lambda = \dfrac{\beta}{\alpha + \beta} \\
                &\iff \beta = \lambda (\alpha + \beta) \iff \beta - \lambda \beta = \lambda \alpha \iff \beta (1 - \lambda) = \lambda \alpha \\
                &\iff \dfrac{\beta}{\alpha} = \dfrac{\lambda}{1 - \lambda} \iff \begin{cases}
                    \alpha = 1 - \lambda \\
                    \beta = \lambda
                \end{cases}
            \end{align*}

            Ainsi,
            \begin{align*}
                \vec{f}(\lambda \vec{u}) &= \vec{f}(\overrightarrow{OP}) \\
                &= \overrightarrow{f(O) f(P)} \\
                &= \overrightarrow{f(O) f(bar((O, 1 - \lambda), (M, \lambda)))} \\
                &= \overrightarrow{f(O) bar((f(O), 1 - \lambda), (f(M), \lambda))} \quad \text{par hypothèse} \\
                &= \dfrac{1}{(1 - \lambda) + \lambda} ((1 - \lambda) \overrightarrow{f(O) f(O)} + \lambda \overrightarrow{f(O) f(M)}) \\
                &= \lambda \overrightarrow{f(O) f(M)} \\
                &= \lambda \vec{f}(\vec{u})
            \end{align*}

            \item[Linéarité.] %%%%%%%%%%%%%%%%%%%%%%%
            Soit $\vec{u}, \vec{v} \in \vec{\E}$, il existe $M, N \in \E$ tels que $\vec{u} = \overrightarrow{OM}$ et $\vec{v} = \overrightarrow{ON}$

            Soit $I \in \E$ le barycentre de $(M, 1), (N, 1)$ (le milieu quoi!).

            ainsi on a,
            $$
            \overrightarrow{OI} = \dfrac{\overrightarrow{OM} + \overrightarrow{ON}}{2} = \dfrac{\vec{u} + \vec{v}}{2} \implies \vec{f}(\overrightarrow{OI}) = \vec{f}\left(\dfrac{\vec{u} + \vec{v}}{2}\right)
            $$
            ainsi,
            \begin{align*}
                \vec{f}(\vec{u} + \vec{v}) &= \vec{f}(2 \overrightarrow{OI}) \\
                &= 2 \vec{f}(\overrightarrow{OI}) \\
                &= 2 \vec{f}\left(\dfrac{\vec{u} + \vec{v}}{2}\right) \\
                &= \vec{f}(\vec{u}) + \vec{f}(\vec{v})
            \end{align*}
        \end{description}
        Donc $\vec{f}$ est linéaire et $f$ est affine.
    \end{description}
\end{demonstration}

\begin{proposition}{}{}
    Soient $\E$ et $\mathcal{F}$ deux espaces affines de dimensions finies $n$.

    Soit $(A_0, \dots, A_n)$ un repère affine/barycentrique de $\E$ et \text{$(B_0, \dots, B_n) \in \mathcal{F}^n$}.
    
    \vspace{0.5cm}
    Alors, il existe une unique application affine $f: \E \to \mathcal{F}$ telle que $f(A_i) = B_i$ pour tout $i \in \llbracket 0, n \rrbracket$ et $f$ est un isomorphisme si et seulement si
    $(B_0, \dots, B_n)$ est un repère affine/barycentrique de $\mathcal{F}$.
\end{proposition}

\begin{demonstration}
    \begin{description}
        \item[Unicité.] 
        Soient $f$ et $g$ deux applications affines telles que $f(A_i) = B_i$ et $g(A_i) = B_i$.
        
        Calculons, pour $M \in \E$,
        \begin{align*}
            f(M) &= f(A_0 + \overrightarrow{A_0 M}) \\
            &= f\left(A_0 + \sum_{i=1}^n \lambda_i \overrightarrow{A_0 A_i}\right) \\
            &= f(A_0) + \sum_{i=1}^n \lambda_i \vec{f}(\overrightarrow{A_0 A_i}) \\
            &= B_0 + \sum_{i=1}^n \lambda_i \overrightarrow{B_0 B_i} \\
            &= g(A_0) + \sum_{i=1}^n \lambda_i \vec{g}(\overrightarrow{A_0 A_i}) \\
            &= g\left(A_0 + \sum_{i=1}^n \lambda_i \overrightarrow{A_0 A_i}\right) \\
            &= g(M)
        \end{align*}
        Donc $f = g$ et l'unicité est montrée.

        \item[Bijectivité.]
        On définit une application affine $g: \mathcal{F} \to \E$ telle que pour tout $i \in \llbracket 0, n \rrbracket$, $g(B_i) = A_i$.
        
        Soit $M \in \E$, on a
        \begin{align*}
            g(f(M)) &= g\left(B_0 + \sum_{i=1}^n \lambda_i \overrightarrow{B_0 B_i}\right) \\
            &= g(B_0) + \sum_{i=1}^n \lambda_i \vec{g}(\overrightarrow{B_0 B_i}) \\
            &= A_0 + \sum_{i=1}^n \lambda_i \overrightarrow{A_0 A_i} \\
            &= M
        \end{align*}
        Ainsi $g \circ f = id_{\E}$.

        De même, pour $N \in \mathcal{F}$, on a
        \begin{align*}
            f(g(N)) &= f\left(A_0 + \sum_{i=1}^n \mu_i \overrightarrow{A_0 A_i}\right) \\
            &= f(A_0) + \sum_{i=1}^n \mu_i \vec{f}(\overrightarrow{A_0 A_i}) \\
            &= B_0 + \sum_{i=1}^n \mu_i \overrightarrow{B_0 B_i} \\
            &= N
        \end{align*}
        Ainsi $f \circ g = id_{\mathcal{F}}$.

        Donc $f$ est un isomorphisme.
\end{description}
\end{demonstration}

\end{document}